\chapter{Contribuições Esperadas e Limites da Proposta}

\section{Contribuições esperadas}

Espera-se produzir uma proposta de ensino introdutório de \textit{machine learning} situada no campo da Computação na Educação Básica. A proposta reúne uma plataforma local, uma sequência didática de oito encontros, instrumentos individuais aplicados em diferentes momentos, critérios de análise e orientações de mediação docente. Esse conjunto permitirá examinar como os estudantes relacionam dados, regras, parâmetros e saídas de um modelo e como explicam seus limites.

A proposta deverá oferecer aos estudantes iniciantes um percurso que não depende de programação prévia nem de formalização matemática avançada. A expectativa é criar condições para formular hipóteses, testar uma alteração e explicar a consequência observada, relacionando representações gráficas e tabelas. O interesse da pesquisa não é afirmar domínio técnico profundo dos modelos, mas identificar compreensões conceituais situadas ao longo da intervenção.

A matriz curricular construída para o projeto permitirá explicitar a relação entre os objetivos de aprendizagem, as atividades e as habilidades de Computação do Ensino Médio relacionadas a fundamentos de inteligência artificial, exploração de modelos simples e análise de dados. Essa articulação procura contribuir para discussões sobre a implementação da Computação na Educação Básica sem alegar que a BNCC prescreve algoritmos específicos ou uma forma única de organização curricular.

A investigação também poderá sistematizar uma experiência de docência adaptável a outros contextos de Ensino Médio, desde que sejam respeitadas condições de infraestrutura, mediação, conhecimento prévio e adequação curricular de cada escola. Plataforma e materiais não serão apresentados como solução pronta ou universal. Poderão, contudo, oferecer subsídios para professores interessados em tratar inteligência artificial e dados de forma menos centrada no uso passivo de ferramentas prontas.

O produto também deverá tornar visível o percurso de planejamento que sustenta sua utilização. Não bastará disponibilizar telas, gráficos e controles. A sequência, os cenários, as perguntas de previsão e os critérios de análise precisarão indicar que compreensão se pretende favorecer em cada atividade e que tipo de explicação poderá ser observado.

Essa documentação poderá oferecer aos professores uma referência para justificar adaptações da proposta, sem padronizar aulas para contextos distintos. Nessas adaptações, será necessário preservar a relação entre conteúdo de Computação, Matemática mobilizada e mediação em sala de aula, para que a proposta não seja reduzida a uma demonstração de ferramentas.

\section{Limites e riscos reconhecidos}

O número de encontros e a duração prevista da intervenção limitam a profundidade com que cada modelo poderá ser abordado. Trabalhar cinco modelos introdutórios em oito aulas exige retomadas e comparações constantes, mas não permite supor que uma única exploração resulte em aprendizagem consolidada. O projeto deverá priorizar conceitos e relações fundamentais, deixando de fora formalizações e procedimentos técnicos que ultrapassem o nível introdutório definido.

Diferenças no repertório matemático, na experiência de uso de computador e no acesso a equipamentos podem afetar a participação dos estudantes. A sequência prevê estratégias de acesso e mediação. Essas medidas não eliminam desigualdades de acesso ou de experiência. Os registros docentes ajudarão a interpretar como tais condições interferem nas produções individuais, evitando atribuir ao modelo ou à plataforma explicação exclusiva para as respostas observadas.

O estudo será realizado em uma turma sob docência do pesquisador, ainda a ser definida. Por esse motivo, suas evidências estarão vinculadas a um contexto específico, à mediação realizada, aos dados sintéticos selecionados e às condições materiais daquele campo. Não se pretende generalizar resultados para todas as escolas públicas ou medir um efeito causal isolado da plataforma.

A comparação entre o núcleo inicial e o núcleo final permitirá examinar mudanças e permanências nas explicações. O desafio de transferência acrescentará evidências sobre o uso dos conceitos em outra situação. Esses resultados não estabelecerão, por si sós, que a plataforma causou as mudanças observadas.

Há limites próprios dos modelos e das representações adotadas. Os conjuntos de dados serão simplificados para tornar relações visíveis, e os resultados produzidos pela interface não deverão ser confundidos com previsões confiáveis sobre pessoas ou situações reais. As vistas tridimensionais representarão métricas calculadas, não atributos adicionais dos dados.

Os elementos didáticos deverão ser revistos antes da aplicação, evitando reforçar a ideia de que sistemas de inteligência artificial são neutros, autônomos ou infalíveis.
